\documentclass[11pt]{article}
\usepackage[utf8]{inputenc}
\usepackage{amsmath, amssymb, amsthm}
\usepackage{booktabs}
\usepackage{hyperref}
\usepackage{graphicx}
\usepackage{geometry}
\geometry{margin=1in}

\newtheorem{theorem}{Theorem}
\newtheorem{lemma}[theorem]{Lemma}
\newtheorem{definition}[theorem]{Definition}
\newtheorem{proposition}[theorem]{Proposition}

\title{Post-Quantum Finality for Lux Quasar \\ {\large Replacing BLS Aggregation with ML-DSA + Ringtail}}
\author{Zach Kelling \\
\textit{Lux Industries Research}}
\date{2026-04-13}

\begin{document}
\maketitle

\begin{abstract}
The current Quasar consensus certificate combines BLS12-381 aggregation with
Pulsar (Ringtail-derived lattice threshold) and ML-DSA lattice signatures
for ``defense in depth.'' We show
that BLS contributes no post-quantum security to the combined cert and
can be removed without loss of PQ soundness. The simplified Quasar
certificate uses only $\mathsf{ML\text{-}DSA}$ (FIPS 204, Module-LWE/MSIS)
and $\mathsf{Pulsar}$ (Ring-LWE threshold; LP-073) --- two independent
lattice assumptions, each post-quantum. The result is operationally simpler,
smaller in the common case, and provably PQ-safe under either assumption
independently.

\paragraph{Naming.} ``Ringtail'' is the underlying signing primitive
(the two-round MAC-authenticated lattice scheme of \cite{ringtail2025});
``Pulsar'' is the Lux systems-layer composition that wraps it with
key-era lifecycle, verifiable resharing, activation cert, and rollback
suitable for permissionless validator rotation (LP-073 / \texttt{github.com/luxfi/pulsar}).
The soundness arguments below apply to the signing primitive; lifecycle
additions are orthogonal.

This document specifies:
\begin{enumerate}
  \item the reduced three-layer Quasar cert: per-block committee cert,
        per-epoch Pulsar pulse, per-era Z-chain ZK proof;
  \item the security argument (forgery requires both MLWE and RLWE breaks);
  \item concrete efficiency numbers measured on MacBook M3 / arm64;
  \item why a fixed-size committee of $k=128$ is scale-invariant from
        $N=10^3$ to $N=10^5$ validators.
\end{enumerate}
\end{abstract}

\section{Motivation}

The classical BLS12-381 pairing-based aggregate signature is broken in
polynomial time by a sufficiently large quantum computer (Shor's algorithm
on the discrete log in $\mathbb{F}_{q^{12}}$). Including BLS in a
``PQ-safe'' finality cert therefore provides only a classical speed-up;
it cannot contribute to the post-quantum security argument.

This leaves Quasar with three natural reasons to keep BLS:
\begin{enumerate}
  \item \textbf{Speed}: classical verification of a BLS aggregate is ${\sim}$1 ms.
  \item \textbf{Size}: a BLS aggregate is 48 bytes; 128 individual ML-DSA
        signatures are $128 \times 2420 = 309{,}760$ bytes.
  \item \textbf{Hedge}: if all lattice assumptions break simultaneously,
        BLS still holds until a quantum attacker materialises.
\end{enumerate}

Against these we weigh:
\begin{itemize}
  \item \textbf{False sense of security}: operators reading ``PQ cert''
        may assume the cert is PQ when in fact it is only PQ conditional
        on verifying the lattice component.
  \item \textbf{Complexity cost}: three independent crypto stacks
        (pairing, Module-LWE, Ring-LWE) to maintain, benchmark, and audit.
  \item \textbf{PQ-first deployments}: networks that explicitly reject
        any classical primitive (regulated finance, long-archival chains)
        cannot use the BLS-inclusive cert at all.
\end{itemize}

We conclude the hedge argument does not justify the complexity and drop
BLS.

\section{Reduced Quasar Certificate}

\subsection{Per-block committee certificate (hot path)}

For each block at height $h$:
\begin{enumerate}
  \item A committee $C_h \subset \mathcal{V}$ of size $k = 128$ is sampled
        from the validator set $\mathcal{V}$ (of size $N$) via a
        post-quantum VRF:
        \[
          C_h = \textsf{Sample}_{\textsf{VRF}}\!\left(\textsf{prev\_hash} \,\|\, h; \textsf{stake}\right).
        \]
        The VRF uses a lattice-based construction (or $\mathsf{ML\text{-}DSA}$
        itself as a hash-based VRF, with unpredictability in the ROM).
  \item Each $v_i \in C_h$ signs the block header $B_h$ with its individual
        $\mathsf{ML\text{-}DSA}\text{-}65$ key:
        \[
          \sigma_i = \mathsf{ML\text{-}DSA.Sign}(sk_i, B_h).
        \]
  \item The block carries $\textsf{cert}_h = \big(C_h^{\textsf{bitmap}}, (\sigma_i)_{i \in C_h^{\textsf{signed}}}\big)$.
  \item A block is considered \emph{locally final} iff
        $|C_h^{\textsf{signed}}| \geq \lceil 2k/3 \rceil$.
\end{enumerate}

This cert is $\leq 128 \times 2420 = 309{,}760$ bytes and verifies in
$\sim 5$ ms on commodity hardware (measured below).

\subsection{Per-epoch Ringtail aggregate (warm path)}

An \emph{epoch} spans $E$ blocks (target: $E = 128$, one hour at 30-second
block times). At the end of each epoch, the committee that signed block
$h_{\text{end}}$ produces a $\mathsf{Ringtail}$ threshold signature over
the epoch tail Merkle root:
\[
  \sigma_{\text{epoch}} = \mathsf{Ringtail.Combine}\!\Big(\big(\mathsf{Ringtail.Sign}_i(r_{\text{epoch}})\big)_{i \in C_{h_{\text{end}}}^{\textsf{signed}}}\Big),
\]
where $r_{\text{epoch}}$ is the Merkle root over $\{B_h\}_{h \in \text{epoch}}$.
$\sigma_{\text{epoch}}$ is $\sim 5$--$10$ kB regardless of epoch length.

Once an epoch aggregate is available, validators may prune the individual
block certs inside the epoch; archival nodes keep them for auditability.

\subsection{Per-era Z-chain ZK proof (cold path)}

An \emph{era} spans $R$ epochs (target: $R = 168$, one week at hourly epochs).
At the end of each era, a Groth16 or PLONK proof $\pi_{\text{era}}$ is
produced on the Z-chain proving:
\[
  \pi_{\text{era}} : \exists\ (C_h, \sigma_h)_{h \in \text{era}}\ \ \text{s.t.}\ \ \bigwedge_h \mathsf{QuasarVerify}(B_h, C_h, \sigma_h) = 1.
\]
$\pi_{\text{era}}$ is $\leq 200$ bytes and verifies in $\sim 3$ ms.
Light clients sync by downloading era proofs, not block certs.

\section{Security}

\subsection{Threat model}

Let $f$ be the Byzantine fraction of $\mathcal{V}$. Classical Byzantine
thresholds require $f < 1/3$. Let the adversary $\mathcal{A}$ be any
probabilistic polynomial time quantum machine (BQP).

\subsection{Forgery reduction}

\begin{theorem}[PQ soundness of Quasar committee cert]
\label{thm:soundness}
Assume $\mathsf{ML\text{-}DSA}\text{-}65$ is existentially unforgeable under
the $\mathsf{MLWE}_{q,k,\ell,\chi}$ and $\mathsf{MSIS}$ assumptions, and
that the VRF output is pseudorandom in the quantum random oracle model.

Then an adversary $\mathcal{A}$ that produces a valid
$\textsf{cert}_h$ for a block $B_h$ not signed by $\geq 2k/3$ honest
committee members breaks $\mathsf{ML\text{-}DSA}$ or predicts the VRF
with non-negligible advantage.
\end{theorem}

\begin{proof}[Proof sketch]
By a standard hybrid argument:
\begin{itemize}
  \item[\textbf{H0:}] Real game. $\mathcal{A}$ forges $\textsf{cert}_h$
        containing $\geq 2k/3$ valid $\mathsf{ML\text{-}DSA}$ signatures
        over $B_h$ from committee members.
  \item[\textbf{H1:}] Replace VRF output with a uniform random sampling.
        Indistinguishable by VRF PQ-pseudorandomness.
  \item[\textbf{H2:}] Committee is independent of $\mathcal{A}$'s choice.
        Probability of an unchallenged forgery on $\geq 2k/3$ honest
        keys is at most $\binom{k}{2k/3} \cdot \varepsilon_{\text{ML-DSA}}^{2k/3}$
        where $\varepsilon_{\text{ML-DSA}}$ is the per-signature forgery
        advantage. For $k = 128$ and $\varepsilon_{\text{ML-DSA}} \leq 2^{-192}$
        (Level III), this is $\leq 2^{-86} \cdot 2^{-192 \cdot 85} \approx 2^{-16 \cdot 10^3}$.
\end{itemize}
The adversary's overall advantage is bounded by
$\mathsf{Adv}^{\textsf{VRF}} + \binom{k}{2k/3} \varepsilon_{\text{ML-DSA}}^{2k/3}$,
which is negligible under our assumptions.
\end{proof}

\subsection{Independence of the two lattice assumptions}

$\mathsf{ML\text{-}DSA}$ is secure under $\mathsf{MLWE}$ (Module-LWE) and
$\mathsf{MSIS}$ (Module-SIS). $\mathsf{Ringtail}$ is secure under
$\mathsf{RLWE}$ (Ring-LWE) with different ring structure and parameters.

A reduction from $\mathsf{MLWE}$ to $\mathsf{RLWE}$ (or vice versa) at
concrete parameters is unknown. In particular, structural attacks against
one family have not transferred to the other in any published result.

\begin{proposition}[Hedged PQ security]
Under the conjunction ``$\mathsf{MLWE}$ is hard $\lor$
$\mathsf{RLWE}$ is hard,'' the combined per-epoch certificate
$(\textsf{cert}_h \text{ for } h \in \text{epoch}, \sigma_{\text{epoch}})$
is unforgeable.
\end{proposition}

This gives us defense in depth strictly within the post-quantum setting,
without relying on any classical primitive.

\section{Why BLS contributes no PQ security}

\begin{lemma}
A BLS12-381 aggregate signature $\sigma_{\text{BLS}}$ can be forged
efficiently by a quantum adversary running Shor's algorithm on the
discrete log problem in the extension field $\mathbb{F}_{q^{12}}$
underlying the pairing target group.
\end{lemma}

Consequently, any certificate whose verification \emph{requires}
$\sigma_{\text{BLS}}$ to be valid can be broken by breaking BLS alone.
If the Quasar cert verification logic is
$\mathsf{Verify}_{\text{BLS}}(\cdot) \land \mathsf{Verify}_{\text{PQ}}(\cdot)$,
then forging the cert still requires breaking the PQ components, so
BLS adds a verification cost without adding an attack barrier.
If the verification is $\mathsf{Verify}_{\text{BLS}}(\cdot) \lor \mathsf{Verify}_{\text{PQ}}(\cdot)$,
then forging the cert requires breaking \emph{only} BLS, reducing PQ
security to zero.

In practice Quasar uses the $\land$-form, so BLS is safe but superfluous.
Since it does not strengthen the cert, we remove it.

\section{Scale invariance of fixed-committee signing}

\subsection{Claim}

For a fixed committee size $k$, the per-block cert generation and
verification cost is $O(k)$, independent of the total validator count $N$.
Only the VRF sampling is touched by $N$, and the VRF sampling cost
is $O(\log N)$ with a weighted-reservoir or alias-method implementation.

\subsection{Measured costs on MacBook M3 (10-core arm64)}

Using \texttt{mldsa-bench} at $\mathsf{ML\text{-}DSA}\text{-}44$:

\begin{center}
\begin{tabular}{rrrrrr}
\toprule
$N$ & $k$ & Keygen & Sign (parallel) & Verify (parallel) & Cert size \\
\midrule
$10^3$  & 128 & 14.8 ms & 3.8 ms & 4.4 ms & 310 kB \\
$10^4$  & 128 & 14.8 ms & 3.8 ms & 4.4 ms & 310 kB \\
$10^5$  & 128 & 14.8 ms & 3.8 ms & 4.4 ms & 310 kB \\
\bottomrule
\end{tabular}
\end{center}

(Keygen is amortized across validator lifetime and only involves the
committee members for a given block; in practice keys are generated
once at validator registration. Sign and verify scale linearly with $k$
only.)

\subsection{Disk and bandwidth}

\begin{center}
\begin{tabular}{lrrr}
\toprule
 & Per block & Per epoch (128 blocks) & Per era (168 epochs) \\
\midrule
Raw ML-DSA committee cert & 310 kB & 40 MB & 6.7 GB \\
After Ringtail roll-up    & —      & 5--10 kB & 1.2 MB \\
After Z-chain ZK proof    & —      & —         & 200 bytes \\
\bottomrule
\end{tabular}
\end{center}

An archival node retains all raw certs (6.7 GB/era). A full node keeps the
current era's raw certs plus all past Ringtail epoch aggregates (${\sim}200$
MB/year). A light client synchronises from Z-chain era proofs only
(${\sim}10$ kB/year).

\section{Recommendation}

\begin{enumerate}
  \item Adopt the reduced Quasar cert (ML-DSA + Ringtail + Z-chain ZK).
        Remove BLS from Quasar cert generation and verification.
  \item Keep BLS available as an \emph{opt-in} classical speed-up for
        deployments that explicitly accept non-PQ security in exchange
        for sub-millisecond verification of a 96-byte aggregate. Gate
        via a chain-configuration flag, off by default.
  \item Use $k = 128$ fixed committee size for all network sizes in
        the range $10^3 \leq N \leq 10^5$.
  \item Use $\mathsf{ML\text{-}DSA}\text{-}65$ (NIST Level III) for the
        individual committee signatures.
  \item Require a PQ VRF for committee sampling. Either a lattice VRF
        (see Boneh--Dai 2024) or $\mathsf{ML\text{-}DSA}$ output used
        as a pseudorandom source in the QROM.
  \item Re-anchor every epoch via a $\mathsf{Ringtail}$ threshold
        aggregate for space efficiency.
  \item Re-anchor every era via a Z-chain $\mathsf{Groth16}$ or
        $\mathsf{PLONK}$ proof for light-client succinctness.
\end{enumerate}

\section{Open questions}

\begin{itemize}
  \item \textbf{PQ VRF}: which construction best balances cost and
        PQ-security? Lattice VRFs are still an active research area.
  \item \textbf{Stake-weighted sampling with liveness}: how do we
        recover when $>1/3$ of stake is offline in a specific epoch?
        Committee rotation helps but a formal liveness bound is needed.
  \item \textbf{Ringtail parameter selection}: parameters for
        $k = 128$ threshold Ringtail have not been published; we may
        need a sensitivity analysis before adopting.
\end{itemize}

\section{Related work}

\begin{itemize}
  \item Celi, del Pino, Espitau, Niot, Prest.
        \emph{Efficient Threshold ML-DSA.} USENIX Security 2026.
  \item Boschini et al.
        \emph{Ringtail: Practical Two-Round Threshold Signatures from
        Learning With Errors.} IEEE S\&P 2025.
  \item Ethereum beacon chain attester committee design (BLS-based).
  \item Avalanche / Snowman probabilistic sampling family.
\end{itemize}

\end{document}
